\section{Modélisation du problème}

\begin{tcolorbox}[colback=black!5,colframe=black,title=]
Cette section présente la modélisation du problème de navigation sous la forme d'un processus de décision de Markov (MDP). Les différentes composantes du modèle, à savoir l'espace des états, l'espace des actions et la fonction de récompense, sont détaillées dans les sous-sections suivantes.
\end{tcolorbox}

\noindent Le MDP est représenté par le tuple $(S, \; A, \; P, \; R, \; \gamma)$ où :

\vspace{10pt}

\begin{itemize}
    \item $S$ est l’ensemble des états qui vont permettre à l'agent de sélectionner une action,
    \item $A$ est l’ensemble des actions qui correspondent aux décisions de l'agent,
    \item $P : S \times A \times S \rightarrow [0, \; 1]$ est la probabilité de transition entre les états, qui modélise la densité de probabilité du prochain état $\mathbf{s}_{t+1}$, conditionnellement à l'état courant $\mathbf{s}_t$ et l'action $\mathbf{a}_t$,
    \item $R : S \times A \rightarrow [r_\text{min}, \; r_\text{max}]$ est la fonction de récompense qui attribue une valeur à chaque transition,
    \item et $\gamma$ est le facteur d'actualisation qui modélise l'importance du futur.
\end{itemize}

\vspace{10pt}

\noindent \textbf{L'objectif de l'agent est d'atteindre la position cible $(x_g, \; y_g)$ tout en évitant toute collision, en sélectionnant une action à chaque instant $t$.}


%%%%%%%%%%%%%%%%%%%%%%%%%%%%%%%%%%%%%%%%%%%%%%%%%%%%%%%%%%%%%%%%%%%%%%%%%%%%%%%%%%%%%%%
%%%%%%%%%%%%%%%%%%%%%%%%%%%%%%%%%%%%%%%%%%%%%%%%%%%%%%%%%%%%%%%%%%%%%%%%%%%%%%%%%%%%%%%
%%%%%%%%%%%%%%%%%%%%%%%%%%%%%%%%%%%%%%%%%%%%%%%%%%%%%%%%%%%%%%%%%%%%%%%%%%%%%%%%%%%%%%%


\subsection{Espace des états}

\noindent À chaque instant $t$, l'agent observe un état $\mathbf{s}_t \in S$ défini par :
$$
\mathbf{s}_t = (\mathbf{M}_t, \; d_t, \; \mathbf{e}_t, \; \mathbf{v}_t, \; \mathbf{o}_t),
$$
où :
\begin{itemize}

    \item $\mathbf{M}_t \in \{0,1\}^{K \times W \times H}$ représente l'environnement local de l'agent construit à partir des $K$ dernières observations. Chaque observation est une carte encodée de taille $W \times H$, centrée sur l'agent. Une cellule vaut 1 lorsqu'elle est occupée par un obstacle et 0 sinon.

    \item $d_t \in \mathbb{R}^{+}$ représente la distance euclidienne à l'objectif :
    $$
    d_t=\sqrt{(x_g-x_t)^2+(y_g-y_t)^2},
    $$
    où $(x_t, \; y_t)$ désigne la position de l'agent et $(x_g,y_g)$ la
    position de la cible.
    
    \item $\mathbf{e}_t=(\cos(\Delta\theta_t), \; \sin(\Delta\theta_t))$ représente l'erreur d'orientation de l'agent par rapport à l'objectif, avec :
    $$
    \Delta\theta_t=\theta_t^{g}-\theta_t,
    $$
    où $\theta_t^{g}$ désigne la direction de l'objectif dans le repère global.
    
    \item $\mathbf{v}_t = (v_t, \; \omega_t)$ représente respectivement la vitesse linéaire et la vitesse angulaire du robot.
    
    \item $\mathbf{o}_t = (\cos(\theta_t), \; \sin(\theta_t))$ représente l'orientation du robot, où $\theta_t$ correspond à son angle d'orientation dans le repère global.
\end{itemize}

\vspace{10pt}

\noindent Les angles sont représentés à l'aide de leurs composantes trigonométriques afin d'éviter les discontinuités associées à une représentation angulaire directe. L'orientation du robot $\theta_t$, est définie sur l'intervalle $[-\pi, \; \pi]$. Les commandes de contrôle sont contraintes par les limites physiques du robot telles que $v_t \in [v_\text{min}, \; v_\text{max}]$ et $\omega_t \in [\omega_\text{min}, \; \omega_\text{max}]$. L'ensemble des paramètres liés au robot sont récapitulés dans la Table \ref{tab:robot_parameters}.

%%%%%%%%%%%%%%%%%%%%%%%%%%%%%%%%%%%%%%%%%%%%%%%%%%%%%%%%%%%%%%%%%%%%%%%%%%%%%%%%%%%%%%%
%%%%%%%%%%%%%%%%%%%%%%%%%%%%%%%%%%%%%%%%%%%%%%%%%%%%%%%%%%%%%%%%%%%%%%%%%%%%%%%%%%%%%%%
%%%%%%%%%%%%%%%%%%%%%%%%%%%%%%%%%%%%%%%%%%%%%%%%%%%%%%%%%%%%%%%%%%%%%%%%%%%%%%%%%%%%%%%


\subsection{Espace des actions}

\noindent À chaque instant $t$, l'agent choisit une action $\mathbf{a}_t \in A$ qui correspond aux commandes de contrôle appliquées au robot :
$$
\mathbf{a}_t = (v_t, \; \omega_t),
$$
\noindent où $v_t$ et $\omega_t$ désignent respectivement la vitesse linéaire et la vitesse angulaire du robot, avec $v_t \in [v_\text{min}, \; v_\text{max}]$ et $\omega_t \in [\omega_\text{min}, \; \omega_\text{max}]$.

%%%%%%%%%%%%%%%%%%%%%%%%%%%%%%%%%%%%%%%%%%%%%%%%%%%%%%%%%%%%%%%%%%%%%%%%%%%%%%%%%%%%%%%
%%%%%%%%%%%%%%%%%%%%%%%%%%%%%%%%%%%%%%%%%%%%%%%%%%%%%%%%%%%%%%%%%%%%%%%%%%%%%%%%%%%%%%%
%%%%%%%%%%%%%%%%%%%%%%%%%%%%%%%%%%%%%%%%%%%%%%%%%%%%%%%%%%%%%%%%%%%%%%%%%%%%%%%%%%%%%%%


\subsection{Fonction de récompense}

\noindent La fonction de récompense est définie comme la somme de plusieurs termes :
$$
R = R_{\text{g}} + R_{\text{r}} + R_{\text{s}} + R_{\text{c}},
$$
\noindent où $R_{\text{g}}$ récompense la progression vers l'objectif, $R_{\text{r}}$ pénalise les rotations importantes, $R_{\text{s}}$ pénalise la proximité des obstacles et $R_{\text{c}}$ pénalise les collisions.

\vspace{10pt}

\noindent Les notations suivantes sont définies :

\vspace{10pt}

\begin{itemize}
    \item $d^g_t$ désigne la distance entre le robot et la cible à l'instant $t$,

    \item $d^o_t$ désigne la distance entre le robot et l'obstacle le plus proche à l'instant $t$,

    \item $\xi$ représente la distance minimale de sécurité.
\end{itemize}

%%%%%%%%%%%%%%%%%%%%%%%%%%%%%%%%%%%%%%%%%%%%%%%%%%%%%%%%%%%%%%%%%%%%%%%%%%%%%%%%%%%%%%%

\subsubsection{Récompense de progression}

\noindent La récompense de progression est définie par :
$$
    R_\text{g} = \begin{cases}
        \beta_1 \times (d^g_{t-1} - d^g_t) + 2 \times v_t - 1 & \text{si } d^o_t \ge\xi_{\text{s}} \\
        0.2 \times \beta_1 \times (d^g_{t-1} - d^g_t) + 2 \times v_t - 1 & \text{sinon.}
    \end{cases}
    $$
    Un bonus supplémentaire est attribué lorsque l'objectif est atteint (i.e. $d^g_t < \xi_{\text{g}}$) :
    $$
    R_{\text{g}} \leftarrow R_{\text{g}} + r_{\text{goal\_reward}},
    \qquad \text{où }
    r_{\text{goal\_reward}} > 0.
$$

\noindent Cette formulation récompense la diminution de la distance à la cible. Lorsque le robot évolue à proximité d'un obstacle, la récompense de progression est volontairement réduite afin de privilégier la sécurité. Une pénalité constante de $-1$ est appliquée à chaque étape afin d'encourager l'agent à atteindre sa cible en un nombre minimal d'actions. Enfin, un terme proportionnel à la vitesse linéaire favorise des déplacements rapides lorsque ceux-ci sont compatibles avec les autres contraintes.

%%%%%%%%%%%%%%%%%%%%%%%%%%%%%%%%%%%%%%%%%%%%%%%%%%%%%%%%%%%%%%%%%%%%%%%%%%%%%%%%%%%%%%%

\subsubsection{Pénalité de rotation}

\noindent Afin d'éviter les rotations inutiles, une pénalité quadratique est appliquée à la vitesse angulaire :
$$
    R_\text{r} = \beta_2 \times r_{\text{angular\_malus}} \times \omega_t^2,
    \qquad \text{où }
    r_{\text{angular\_malus}} < 0.
$$
\noindent Cette pénalité limite les rotations excessives et favorise des trajectoires plus fluides.

%%%%%%%%%%%%%%%%%%%%%%%%%%%%%%%%%%%%%%%%%%%%%%%%%%%%%%%%%%%%%%%%%%%%%%%%%%%%%%%%%%%%%%%

\subsubsection{Pénalité de sécurité}

\noindent Lorsque la distance au plus proche obstacle devient inférieure au seuil de sécurité $\xi$, une pénalité est appliquée : 
$$
    R_\text{s} = \begin{cases}
        \beta_3 \times r_{\text{safety\_malus}} \times \left( \frac{1}{(d^o_t)^2} - \frac{1}{\xi^2} \right) & \text{si } d^o_t < \xi_{\text{s}}, \\
        0 & \text{sinon},
    \end{cases}
    \qquad \text{où }
    r_{\text{safety\_malus}} < 0.
$$
Cette pénalité augmente rapidement lorsque le robot se rapproche d'un obstacle, ce qui l'incite à conserver une marge de sécurité.

%%%%%%%%%%%%%%%%%%%%%%%%%%%%%%%%%%%%%%%%%%%%%%%%%%%%%%%%%%%%%%%%%%%%%%%%%%%%%%%%%%%%%%%

\subsubsection{Pénalité de collision}

\noindent Enfin, une forte pénalité est attribuée en cas de collision :
$$
    R_{\text{c}} = \begin{cases}
        r_{\text{collision\_malus}} & \text{si } d^o_t < \xi_{\text{c}}, \\
        0 & \text{sinon},
    \end{cases}
    \qquad \text{où }
    r_{\text{collision\_malus}} < 0.
$$
Cette pénalité vise à décourager fortement toute trajectoire conduisant à une collision.

\vspace{10pt}

\noindent Les coefficients de cette fonction de récompense ont été déterminés empiriquement au cours des différentes phases d'expérimentation, en s'appuyant sur les principes proposés dans la littérature \citep{long2018, choi2021, huang2024}. Ils sont récapitulés dans la Table \ref{tab:reward_parameters}. \textbf{L'objectif est d'obtenir un compromis entre rapidité de déplacement, fluidité des trajectoires et sécurité vis-à-vis des obstacles.}
